\subsection{An unramified height-two obstruction beyond the first level}

\begin{proposition}[B-RD22: four-torsion does not suffice at $p=2,e=1$]
\label{prop:B-RD22}
\status{Proved; exact finite certificate and independent parent and obstructions/crystals audits passed.}
For $p=2$, $e=1$, and every height $h\geq2$, the least index in
B.4 is at least two. In the manuscript's indexing, even
$M(1)=M[4]$ does not determine the character on inertia elements of
odd order. This already occurs over the unramified quadratic field
$K=\mathbf Q_2(\zeta_3)$.
\end{proposition}
\begin{proof}
Let
\[
 E/K:\quad y^2=x(x-1)(x+3),\qquad
 M=E[2^\infty],\qquad
 N=\mu_{2^\infty}\oplus\mathbf Q_2/\mathbf Z_2.
\]
The group $N$ extends to a split Barsotti--Tate group over
$\mathcal O_K$. Put $a=1+2\zeta_3$, so that $a^2=-3$, and choose
$i^2=-1$. The points
\[
       P=(a,ia(a-1)),\qquad Q=(3,6)
\]
satisfy $2P=(0,0)$ and $2Q=(1,0)$, and hence form a basis of
$E[4]$. Every element of $G_K$ fixes $a$ and sends $i$ to either
$i$ or $-i$. Consequently it sends $P$ to $\chi_2 P$ modulo four
and fixes $Q$. The equivalence between finite etale group schemes
over the characteristic-zero field $K$ and finite $G_K$-sets gives
an isomorphism of group schemes
\[
                     M[4]\simeq\mu_4\oplus\mathbf Z/4
                              =N[4].                 \tag{B-RD22.1}
\]

For the displayed integral Weierstrass equation the exact invariants
are
\[
 \Delta=2^8\,3^2,\qquad c_4=2^4\,13,\qquad
 j=\frac{2^4\,13^3}{9}.
\]
The discriminant valuation is eight, which is minimal since every
integral change to another Weierstrass equation changes it by a
multiple of twelve and a discriminant of an integral equation has
nonnegative valuation. In particular the curve is bad. Its integral
$j$ gives potential good reduction
\cite[Section~2, p.496]{SerreTate1968Good}. Since $v_K(j)=4>0$,
the good special fibre has $j=0$. In characteristic two its equation
has the form $y^2+a_3y=x^3+a_4x+a_6$, with $a_3\ne0$
\cite[Section~9.1, p.8]{AnemaTopTuijp2018Hesse}. Negation sends
$(x,y)$ to $(x,y+a_3)$, so there is no nonzero geometric
two-torsion point. Thus the good fibre is supersingular.

Pass to the completed maximal unramified extension and let $L/K$
be a minimal good-reduction extension, with inertia image of order
$m$. A change to a unit-discriminant equation over $L$ shows that
\[
                            12\mid8m,
\]
so $3\mid m$. The inertia image therefore contains an automorphism
of order three. On the supersingular special elliptic curve its
action on the rational Dieudonne module has eigenvalues $\zeta_3$
and $\zeta_3^{-1}$, hence trace $-1$: equivalently, use the standard
supersingular equation $y^2+y=x^3$ and its order-three automorphism
$(x,y)\mapsto(\zeta_3x,y)$; the two eigenvalues are interchanged by
Frobenius. The order-three trace of $N$, which has good reduction,
is $2$. Thus (B-RD22.1) does not determine the requested character.
Both groups satisfy the manuscript's monodromy condition.

For height $h>2$, add $h-2$ copies of the constant etale group
$\mathbf Q_2/\mathbf Z_2$ to both sides. Their equal traces add the
same constant, so the discrepancy remains. The entire finite
four-torsion calculation is checked with exact rational arithmetic
in $\mathbf Q(\sqrt{-3},i)$ by
\path{code/verify_p2_character_level2.py}, including all sixteen
points, all $256$ additions, and the Galois action.
\end{proof}
